\documentclass{article}

\input{properties/packages.tex}

\input{properties/aliases}

\input{properties/settings.tex}


\begin{document}

\begin{center}{\Large \textbf{
Reading notes
}}\end{center}

\begin{center}
Christian Boyd
\end{center}

\begin{center}
{\small \sf Boyd.Christian.M@gmail.com}
\end{center}

\begin{center}
\today
\end{center}


\section*{Abstract}
{\bf
These are my reading notes, mostly focused on finance and/or financial planning.
}

\vspace{10pt}
\noindent\rule{\textwidth}{1pt}
\tableofcontents
\noindent\rule{\textwidth}{1pt}
\vspace{10pt}


\section{Money creation in the modern economy (M. McLeay, A. Radia, \& R. Thomas)}

After reading \cite{BOE2014}, I felt that there wasn't much substance to the actual article; however, it has led to some other references on money supply and inflation.


\section{Will the securitization revolution spread? (R. Feldman)}

In \cite{Feldman1995}, Feldman considers why small business loans have not been securitized to the extent of mortgages, auto loans, and credit card receivables.  This discussion, however, is from 1995 and so this should be more of a jumping-off point than taken as an explanation of the modern situation for small business loan securitization.

Feldman notes that the benefits to securitization are obvious (increased liquidity for lenders, increased loans for borrowers, stable income opportunities for investors), so a natural suspect for the lack of adoption is the difficulty (i.e., cost) in securitization.
\begin{itemize}
    \item Small business loans require detailed information on the management team and the local economy in which the business operates.
    \item Small business loans lack (or, at least they might've in the past) standardized terms.
    \item There is a lack of publicly-available, long-term performance data on small business loans.  I.e., it's a more opaque market when compared to mortgage loans.
\end{itemize}

I find this first point interesting, as it should presumably apply for mortgages and credit card receivables.  I.e., a bank requires detailed income and credit history for a mortgage applicant and a credit card company needs the same to establish a revolving line of credit.  \note{[Consider]:} {\it how does the future income calculation differ between a small business and an individual?}  Feldman notes that this cost/lack of understanding is mostly on the investor side (lenders are already capable of making the loans), who will choose a better-understood investment like mortgage-backed securities when given the option.

Feldman further notes that investors want loss protections when they don't understand the mechanics of the securities they're purchasing.  These loss protections, however, imply that the bank (or, more generally, the lender) still carries the credit risk of the securitized loans.  Maintainin the credit risk requires banks to buy costly insurance against losses on securitized loans.  \note{[Consider]:} {\it how do SBA 7(a) guarantees help with this?}  Moreover, Feldman states that bank regulations consider loans that still carry credit risk to have not been fully sold; i.e., the bank must still meet capital requirements associated with having the securitized loans on its books.  From the bank's point of view, securitization provides liquidity but does not remove the accounting complexity or credit risk associated with the loans.  From the investors point of view, small business loan securities are a complicated--and risky--asset that are not worth the trouble once all the required protections are priced into the return.  It's re-emphasized throughout that the informational burden on the investor is the primary cost preventing broader adoption of securitized small business loans.

\note{[Competition]:} Feldman notes that small business loans are less subject to non-local (lending) competition.  Worth looking into whether this is still the case 30 years later.

\note{[Loan standards]:} it's noted that not all loans are made through the SBA program.  It's worth looking into the size of this non-SBA loan market.


\subsection{Summary}

This has been a thought-provoking article that fleshes out the pros and cons associated with small business loan securitization.  Feldman's points here should be contrasted with the current small business loan market.


\section{R. Shiller - From efficient markets theory to behavioral finance}

In \cite{Shiller2003}, Shiller puts forth that stock price volatility is at odds with the efficient market hypothesis, which he claims requires stock prices to equal the optimal forecast of the present value of all future dividends from that stock.  In particular, Shiller requires market movements in an ``efficient" market to result from new information that causes market participants to rebalance their portfolios.  Consistent volatility that exceeds informational shortcomings, however, would discredit Shiller's construction of the efficient market hypothesis.

Shiller goes on to compare plots of present value calculations against stock price trends over the past century.  He claims that there is too much volatility in the S\&P to positively support the efficient market hypothesis.  Confusingly, it seems that there is some predictive power in the price of individual stocks about their present value, but that no simple weighting or time-discounting scheme appears to reproduce stock prices.  In particular, no simple scheme of determining the present value of stock divident payouts seems to accurately track boom and bust cycles.

Instead, Shiller transitions to discussing behavioral psychology, where people are shown to exhibit irrational strategies (or heuristics) that would lend themselves to boom-bust cycles in trading markets.  He relates the fervor of a Ponzi scheme to market bubbles, where initial investors get such high returns that new investors are no longer considered with the fundamentals that created the returns.  People simply want in on the action and embrace the risk.

Shiller goes on to discuss the efficient market answer to irrational market participants: that ``smart money'' will move opposite the irrational participants to profit off their mistakes.  An issue Shiller brings up is an apparently-widespread issue that it can be difficult to exert downward pressure on an asset bubble.  He notes the example of 3Com's Palm offering, where Palm shares rose to such heights that the remaining 3Com business had an implied negative value on their overall holdings.  An issue in this bubble was that investors who realized the over-pricing of Palm could not easily short Palm stock due to shortages and high price barriers to put options.  While it's stated that this state of affairs is relatively rare, this sounds a lot like discussion around the housing bubble of the mid-00s.  Even though people had concerns about housing prices, there was no obvious/efficient mechanism for skeptical investors to push prices down.\footnote{Of course, there is the story of {\it The Big Short}, where clever investors made money through credit default swaps with unsuspecting banks.  These investors, however, did not have any actual capacity to sell the underlying homes (and push prices down).}

This paper wraps up with a summary that the stock market certainly reflects information about the companies being bought and sold.  It may be surprising in how efficient it ``normally'' operates, but it is certainly not perfectly efficient.  Shiller advises that economists (or investors, or anyone really) should keep in mind that the market prices are not perfect and that irrational behaviors are a necessary consequence of humans being the buyers and sellers in the market.  I think this paper was most interesting to me as an introduction to how Shiller--and others--value stocks, and think of them as dividend return vehicles.  Any alternative pricing strategy appears to be implicitly ``irrational.''


\note{[Consider]}:
\begin{itemize}
    \item Do I agree that the error is necessarily uncorrelated with the current price/information?
    \item E.g., if a crucial piece of information is missing from the pubilc knowledge, and that piece of information would drastically alter the price, wouldn't the error necessarily depend on the price?
    \item Cast another way: if the error is the actual minus the forecast, then how can this residual not depend on the forecast?
    \item What {\it causes} the uncertainty in $P_t^*$?  Is it purely the correct discount rate for each future dividend?  Presumably it is based off the actual dividends that are eventuallly paid out, rather than being another estimate.
    \item In what universe is the uncertainty in the present value of all future {\it known} dividend payouts bounded (below) by the uncertainty in the market forecast of that present value?  If the fed guaranteed 100 years of interest rates in quarterly intervals, how would the uncertainty in the forecasted present value of future possible dividend payments continue to bound the uncertainty in the present value of actual (known) future dividend payments?
    \item Am I misunderstanding the concept of {\it variance} in this discussion?  Why would the forecast of a random variable ever have a smaller variance (i.e., uncertainty) than the variable itself?
    \item How can it be that per-stock analyses show reasonable price forecasting through the present value, but that this apparently fails on the broader market?  Is this a tail effect where the ``market average'' is skewed by the exceptionally unpredictable stocks?  Do portfolios of stocks somehow create more uncertain forecasts than individual ones?
    \item If boom and bust cycles are removed ({\it if} they can be removed), does the stock price more accuratly reflect present value?  Are boom/bust cycles an aggregate market phenomenon that detracts from an otherwise present-value forecasted stock pricing market?
\end{itemize}


\section{K. Cole, J. Helwege, and D. Laster - Stock market valuation indicators: is this time different?}

This paper discusses the state of the stock market following the summer of 1995.  It's noted that the 2.5\% dividend yield in the summer of 1995 was a historical low.  It's further noted that the market-to-book ratio over 3 exceeds the historical average of 1.9.  It starts out by noting that low dividend yields and high market-to-book ratios (stock market capitalization relative to the accounting book value of a company) tend to signal poor future performance.  Counter-arguments to this thinking is allegedly that the era of tax-advantaged stock buybacks and new accounting rules have resulted in an artificial increase in market-to-book ratios and an artificial decrease in dividend yields.  {\it Note: I can't help but notice the echoes of Robert Shiller's critique of the efficient market hypothesis, where he mentioned that new techniques and rationales emerge in a bubble to support ``the new norm'' of eternally-high stock prices and metrics.}  The trio aim to rebuke these claims and sound the alarm on over-priced stocks in this paper.

There's some interesting discussion about share repurchases vs. dividend payouts.  Aside from the obvious tradeoffs (taxes, primarily), it's noted that a lot of the negative connotation with decreased dividend yields or canceled dividends doesn't seem to weigh down share repurchases.  For example, announced share buybacks that don't come to fruition or scaled back share buyback plans don't seem to cause investor alarm.  From this perpspective, the trio consider whether share repurchases are equally valuable to investors, on a dollar-per-dollar basis.  It's noted further that share issuances acts to negate share buybacks.  On this basis, the net share buybacks (share repurchases minus share issuances) is claimed to still highlight low investor return in the 90s despite high stock prices.

The paper concludes by more or less summarizing that, even when share buybacks and new tax rules are taken into account, the current (1995) market appeared over-valued.  Of course, the stock market (as measured via the S\&P 500) would continue roaring until late 2001, where it faultered multiple times until appearing to regain its rapid growth in the 2010s.  The stock market may have been overvalued, but instead of returning to lower level, it repeatedly crashed in the 2000s while trying to grow faster than the fundamentals would support.  Personally, this begs the question of whether market fundamentals ever truly act to drag the price down, or if the market grows as fast as investors will support.

\note{[Consider]}:
\begin{itemize}
    \item The S\&P {\bf 500} (note, not the S\&P 400 used for market cap analysis in the text) dividend yield was a measly 1.17\% in September 2025, an apparent historical low.  I'm guessing from the data that the norm is closer to 2\%, which is still low in the eyes of the text.  A Google search suggests a market-to-book ratio somewhere between 2.5 and 2.9, which is on the higher side but not as inflated as the text describes.  
\end{itemize}


\section{K. Case and R. Shiller - Is there a bubble in the housing market?}

\note{[Todo]}: Read \cite{Case2003}


\section{L. Zeng, P. Luk - Examining share repurchases and the S\&P buyback indices in the U.S. market}

In \cite{Zeng2020}, Liyu Zeng backtests the S\&P 500 Buyback Index against their other indices.

\note{[Caution]}: The 20 year period in question is likely one of the worst for the S\&P 500, as it includes both the dot com bubble and the 2008 financial crisis, with only a decade surrounding these shortcomings.  The annualized rate of return is only 6.1\% over this period, which suggests to me that it is not representative of the long-term 10\% annualized returns typically seen over 30+ year periods.  The buyback index in question was also mostly backtested data, since the index didn't launch until 2013.  These findings should be re-done on a longer time horizon, possibly in 2030 if the era of buybacks is considered to be post-2000.

The S\&P 500 Buyback Index tracks the 100 companies in the S\&P 500 with the highest buyback ratio, where the buyback ratio is the trailing 12 month amount of share repurchases divided by the total market capitalization at the {\it beginning} of the trailing 12 month period.  Curiously, this is an {\it equally-weighted} index.  This index is apparently adjusted quarterly in January, April, July, and October.

The returns of the S\&P 500 Buyback Index are analyzed over the 20 year period from 1999 to 2019.  The buyback index apparently outperformed the standard S\&P 500 over most of the 20 years, and the buyback index is supposed to have experienced 5.5\% higher returns per year in an annualized basis.

\note{[Yield is king]}: it's not emphasized, but the yield portfolio--the top 100 companies in the S\&P 500 when both dividends and buybacks are used in the trailing 12 month calculation--out-performed the buyback index with lower volatility.  There's an aside that the yield index didn't out-perform the S\&P 500 in as many years as the buyback index, but I'm not sure why the yearly comparison would matter more than the annualized returns over the 20 year period.

Exhibit 6 is an interesting set of data.  Of the five indices compared (buyback, dividend, all-yield, S\&P 500, and {\it equally-weighted} S\&P 500), the standard S\&P 500 performed the worst over the 20 year period.  I would be very curious to see 30 year data, if it's available.  If anything, I think this chart is saying that recent-term yield is a better indicator of future returns than current (or recent) market capitalization.  That feels a bit tautological, but perhaps it's interesting that a one-year run up is enough time to gauge a company's value-returning habits.  The equally-weighted S\&P 500 out-performing the standard index over the 20 year period suggests that there are some large capitalization companies not returning value to shareholders as effectively as other smaller--but still large--capitalization companies.

There's a discussion that the bull cycle momentum catapults the buyback index over the standard index in up months vs. down months.  Exhibit 9 mostly demonstrates that dividends anchor losses against down cycles, and capital gains anchor gains in growth cycles.  The buyback index has the largest upside in growth cycles, but the all-yield index has much better protection against down cycles with similar upside in growth cycles.  Again, this seems to imply that the trailing 12-month shareholder returns are a more valuable weighting metric than purely the market capitalization.  That would be an interesting index to try and maintain\ldots

The discussion of the small-cap weighting of these yield indices (buyback, dividend, and all-yield) really begs for a modern analysis, where the past 5-10 years of the S\&P 500 have been dominated by the top 6-7 tech stocks.  I wonder if the small-cap bias would invert the relation from exhibit 9, where the S\&P 500 was out-perfomed by the other indices that weighted indifferently to market capitalization.

In the small- and mid-cap indices, the buyback index really did outperform the all-yield index.  Notaby, the trend of the equal-weight index outperforming the standard S\&P index also held true.  The small- and mid-cap indices actually seemed to better exemplify the analysis of the S\&P 500 buyback index, where growth-cycles propel the small- and mid-cap buyback indices past the standard index and down cycles don't cause as much value loss.  The all-yield indices both showed negative performance in down cycles vs. the standard (weighted by market capitalization) index, where the S\&P 500 all-yield index still maintained increased returns.  The broader trend nevertheless holds, where dividends anchor returns in down cycles and capital gains propel returns in growth cycles.


\section{E. Fama, K. French - A five-factor asset pricing model}

E. Fama and K. French start \cite{Fama2015} with a bit of a history lesson on how people have thought of CAPM and pricing models since the 60s and 90s.  

When introducing the models, it began to sound like the models under discussion were simply adding sufficiently many parameters to capture a majority of the effect.  Funnily enough, this appears to be Fama's preferred interpretation.  They reference \cite{Fama1996} as motivating the intepretation that the portfolios as prescribed are sufficiently distinct and numerous as to span some true, unknown set of market variables that effectively predict the future return.

{\bf Takeaway}:  ``There are patterns in average returns related to Size,
B/M, profitability, and investment.'' I think that quote summarizes the findings, particularly when understood as a response to CAPM, where portfolio diversification/risk was the primary driver of returns.  This is a very long, wordy article with lots of tables and appendices, but I'm not sure if more is to be gleaned there.  My understanding of the article is that there are a handful of market factors driving the majority of returns (i.e., rather than there being an innumerable amount), and that the returns of most portfolios can be understood as arising from a decomposition of that portfolio into these factors.  Whether that's true or not is another matter, but that's how I understood the main thesis of their work.


\section{Wikipedia article on the Capital Asset Pricing Model (CAPM)}

To get a baseline for historical discussions of risk vs. return and portfolio pricing, I've been reading through the Capital Asset Pricing Model (CAPM) wikipedia article \cite{CAPM2026}.  My current understanding of the thesis behind CAPM is that the returns of a given portfolio are essentially determined by the market returns for the asset classes in the portfolio, weighted by their risk.\footnote{I'm not saying this is true or correct, but that it's my impression of the CAPM assumptions.}  The subsequent work by Fama and French \cite{Fama2015} (see their earlier work as well) notably provides evidence for contributions outside of CAPM.  Criticisms in the Wikipedia artcle seem to boil down to the lack of a time-constant $\beta$, the immeasurability of {\it expected} market returns for all market participants, and the assumption that the risk profile is totally captured by the mean and variance (e.g., ignoring asymmetric downside risk). 

I won't go much further fleshing out the Wikipedia page, but want to keep the link \cite{CAPM2026} for its references (both within and outside of Wikipedia).


\section{Wikipedia article on the Arbitrage Pricing Theory (APT)}

One of the links that piqued my interest on Wikipedia's CAPM page \cite{CAPM2026} is the Arbitrage Pricing Model \cite{APT2026}.  The APT is apparently based on the {\it Law Of One Price} (LOOP), where identical goods sold in different markets (e.g., locations, vendors, etc.) should be priced identically when expressed in the same currency.  The LOOP only holds insofar as taxes, regulations, transaction costs, and all other trade frictions or market manipulations are absent, as it represents the transactional ideal where all arbitrage opportunities have reach a single price equilibrium.

The Arbitrage Pricing Theory (APT) itself, seems like a generic conceptual framework for how returns ought to be related to intrinsic market factors, and {\it risk premiums} associated with those factors.  I'm having some difficulty making sense of the various representations in the article \cite{APT2026}, but will make my best-effort attempt to summarize my understanding in what follows.  Both the Wikipedi article and the original paper \cite{Ross1976} emphasize the multi-dimensional nature of a portfolio of assets, but I'm currently struggling to see the reason and so will start by considering a single asset whose rate of return $r$ is under study.  The APT states that, as far as returns as concerned, the rate of return can be understood as arising from the risk-free rate $r_f$, a number of {\it factors} $f_i$, and an error term $\epsilon$,
\begin{equation}
    r = r_f + \epsilon + \sum_i \beta_i f_i \quad.
\end{equation}
The rate of return $r$, the error $\epsilon$, and the factors $f_i$ are random variables.  The risk-free rate $r_f$ and $\beta_i$ are merely coefficients, where the $\beta_i$ are $\beta$-like in that they denote the sensitivity of the asset to the $i$th factor, $f_i$.  The major consequence of this framework is that the expected return is given by the risk free rate $r_f$, the $\beta$-like $\beta_i$ coefficients, and the {\it risk premiums} $\mu_i$ associated with the factors $f_i$.  I.e.,
\begin{equation}
    \langle r \rangle = r_f + \sum_i \beta_i \mu_i \quad ,
\end{equation}
where we presume $\langle \epsilon \rangle = 0$ for the error term.

The main distinction of the APT relative to the CAPM that I currently see is that it relaxes the prescriptive power of CAPM to accommodate a multi-factor explanation for asset prices.  My current understanding is that this kind of framework underlies the modern multi-factor models like the 5-factor model \cite{Fama2015}.  For example, the framework doesn't say what the underlying factors are, aside from suggestively naming the resulting coefficients {\it risk premiums}.  All that we know is there are a linear combination of factors, which increase the return of an asset relative to a reference risk-free rate.  


\section{S. Ross - The arbitrage theory of capital asset pricing}

\note{[Todo]}: read \cite{Ross1976} after above\ldots  See also the working paper \cite{Ross1971}, where the rigorous framework was apparently fleshed out.


\section{M. Sammon and J. Shim - Index rebalancing and stock market composition}

In \cite{Sammon2026}, the authors consider the interesting case that index funds which blindly track the market necessarily {\it time} the market through rebalanacing.  Naturally, they suspect that these timings are unfavorable, as shares are presumably issued when prices are high and repurchased when prices are low.  The authors test this hypothesis in the context of an relationship they claim between a value-based market portfolio and owning an equal percentage of shares outstanding in all companies, or an ownership ratio.  Through this relationship, they state that index funds (or indexes themselves) must rebalance to maintain their desired ownership ratio.

\subsection{Notes on the Introduction}

They introduce the concept of {\it Market Tracking Error} (MTE), or the inability of popular value-weighted indexes\footnote{I have always said/though ``indices,'' but it seems like the preferred nomenclature is ``indexes.''} to track the actual rapid day-to-day change in the stock market.  The two claim there is a 75-300 bps (or, 0.75\% - 3\%) MTE, which is much greater than the index tracking error (index funds tracking the indexes) stated as 5 bps (0.05\%) or less.  They further claim that, beyond MTE, there is a real performance drag of 46-69 bps (0.46\%-0.69\%) in annualized terms.  They state that the improved returns are dominated by using delayed shares outstanding data (diminishing returns kicking in around 2 years of delayed data) vs. rebalancing less frequently, which is fascinating.

The authors introduce the {\it Market Information Ratio} (MIR), which they define as an index's return minus the daily-rebalanced market return divided by the MTE.  They further claim that delaying updating the shares outstanding data by 1-2 years appears to maximize the MIR.

\note{[On buybacks and issuance]}: there are a lot of references in section 1.1 on related literature to papers on buybacks/issuance as containing predictive power for future returns.  These would be good papers to look into.

It's worth noting that it seems like they do a lot of data cleaning.  This is probably legitimate but they ``remove outliers,'' which may have some impact on the results.  They frequently reference CSRP (Center for Research in Security Prices), Thomson Reuters, and S\&P for data.

\note{[Thoughts on implications of the index lag]}: does the index drag relative to using delayed shares outstanding relate more to market inefficiency, bad timing by management, or simply a long response time for the market to correct the price based on the updated share count?  For example, is this increased performance for delayed values a representation that companies tend to issue and repurchase stock at inopportune times, where buybacks precede underperformance and issuance overperformance?  Or do investors (broadly) have a tough time influencing the stock price due to issuance/repurchase, where additional buy/sell signals due to the open-market transactions aren't sufficient to move the needle?  \note{[Note]}: this seems to be discussed in section 4.

\subsection{Notes on Conceptual Framework}

Their ownership ratio definition seems quite sensible.  Roughly, if an index is meant to weight their portfolio by market capitalization, then they are necessarily sensitive to changes in shares outstanding as they as part of the definition of the market capitalization.  I think the necessity of the ownership ratio is even more important when taking into account that indexes often track the market capitalization implied by the openly-traded shares, or a stock's {\it float}.  This value is even more susceptible to fluctuations during open-mark repurchases or issuance.

\subsection{Notes on the portfolio analysis}

There is an underperformance of roughly -4\% for the portfolio that isolates the amount of the fund dedicated to rebalancing.  The authors suggest that roughly 10\%-15\% of a total fund is involved in rebalancing for any given period (e.g., a quarter), so that this amounts to the roughly 40 bps, or 0.4\% under-performance of value-weighted index funds due to timing the market.

\subsection{Notes on designing a better index}

The authors make an interesting point that index funds settling transactions after closing leads to an invisible trading cost due to indexes using the same timing.  Their point is that both the index and the fund are mistiming the market relative to true "market" returns, which are idealized as settling instantaneously throughout the day.

They find that rebalancing frequency appears less important than using stale shares outstanding data.  Interestingly, there's not a huge impact from rebalancing monthly, quarterly, yearly, or bi-yearly, so long as the shares outstanding data is around 2 years old ($\sim$500 trading days).  This begs the question of whether the market has a hard time pricing due to share issuance/buybacks, or if the corporate actions lead to negative near-term performance.

Curiously, trading costs increase with delayed shares outstanding data, though only marginally up to the roughly 2 year mark (500 trading days) often found to be beneficial for returns.

It seems that the increased returns of using stale data comes largely from market timing, and that using stale data always increases returns in the absence of transaction costs.

\note{[Note]}: do these findings suggest that buying and holding (forever) is superior to maintaining a ``market-like'' portfolio?  I.e., would a fund do better to invest new funds without ever rebalancing prior positions?  Or is using old shares outstanding an improve strategy even when starting from scratch?


\section{D. Ikenberry, et al - Market underreaction to open market share repurchases}

Read \cite{Ikenberry1995}\ldots


\section{C. Asness - Why not 100\% equities?}

The paper \cite{Asness1996} caught my attention as I was watching a podcast discussing a recent paper finding that 100\% equity portfolios out-performed pretty much all combinations of stocks, bonds, and T-bills for retirement performance.\footnote{The actual article in \cite{Asness1996}--rather than the blurb--can be downloaded from a link on the right-hand side.}  I have not fully read this paper (or the inciting 100\% equities paper) through, though the main point seems to be that similar risk and performance to the 100\% equity portfolio can be obtained through leverage (i.e., borrowing) of a more balanced (i.e., diversified) portfolio.

The much more to-the-point online post \cite{Asness2024} essentially summarizes this point.  My understanding from this pointed article, is that stocks {\it should} naturally out-perform bonds over long investment horizons.  The counter-point to an all-equity portfolio is the {\it claim} that a 100\% equity portfolio does not maximize return per risk, where ``risk'' itself becomes tricky to define in the context of long-term vs. short-term investment periods.  In \cite{Asness1996}, it was claimed that a levered 60\% stock/40\% bond investment portfolio actually out-performed an all equity portfolio, once levered to the same ``risk'' of the all-equity portfolio.

I'm not sure if the strongly-worded retort really added much quantitatively, but it is interesting to ask {\it why is leverage viewed with more scrutiny than the ``risk'' it carries?}  If an all-equity portfolio is really as ``risky'' as a leveraged, diversied portfolio, why is the latter more intimidating?  Is it simply a matter of accessibility for the average person trying to save for retirement?  Alternatively, what does a leveraged 401(k) portfolio look like in practice -- how is interest paid, and how are rates set?  Or does an investor simply invest in a leveraged fund, which borrows on behalf of its investors?  Maybe most vexing of all: how is ``risk'' defined -- is there a basis for using a Sharpe ratio (or similar) when looking at long-term investment horizons?


\section{David Blanchett and Jeremy Stempien  - Investment Horizon, Serial Correlation, and Better (Retirement) Portfolios}
In \cite{Blanchett2024}, the authors consider alternative measures of risk for long-term portfolios.  \note{[Note]}: it appears that ``equities" refer to domestic stocks in each country involved in the study, rather than an international holding of equities.  They review data, purportedly demonstrating significant auto-correlation in stocks over long periods; i.e., implying that stocks are less risky on the long-term than their 1-year volatility implies, due to mean-reversion (highs following lows, lows following highs).  They also show an increasing correlation of all examined asset classes (T-bills, bonds, equities, and commodoties) with inflation over longer periods, spanning 1-10 years.  For example, equities appear to demonstrate small negative correlation with inflation over 1 year periods, and increase to smalll positive correlation for 3-10 year periods (increasing with period).

Exhibit 3 is particularly interestin as it visually depicts the aforementioned autocorrelation.  Bonds and bills increase in standard deviation when compared over multi-year periods, and equities and commodities decrease in volatility.  It seems reasonable that, were sequential years uncorrelated, there would be no variation in multi-year period, or at least no clear trends.

Exhibit 5 demonstrates that real historical periods result in the optimal portfolio skewing heavily--if not entirely--toward equities as investment periods increase from 1-20 years.  They also compare to bootstrapped (randomized year data) periods, which more or less lack the dependence on the period.

In Exhibit 6, the main takeaway is the rapid growth of the equity component in the optimal portfolio for even a risk-averse utility function.  When analyzing real returns, the ``high risk aversion'' portfolio goes from around 30\% equities for 1-year investment periods to nearly 60\% for 17-20 year periods.  Naturally, there is little pattern in the plots using bootstrapped returns.  This is to say that--at least according to the chart--even risk-averse portfolios would prefer majority equities over longer (17-20+ year) investment periods.

Exhibit 7 is a curious study analysis where the geometric annual returns are artificially suppressed to those of bonds in the country being studied.  The intent is to remove the equity premium, so that we're only considering the risk profile of equities due to their autocorrelation over long periods compared to bonds.  Even without the equity premium, there is a marked increase in equity allocation in the optimal portfolio looking at real returns, from around 10\% for 1-year periods to nearly 40\% for 20-year periods.  Because the geometric returns are suppressed, this additional preference for equities is claimed to be only in their risk reduction due to the longer-period mean reversion of equities.  

My takeaway here is a trend from related papers that long-term returns really need to be understood through metrics defined over equally long periods.  For example, there's a recurring trend that stocks appear less risky over the long-term in terms of expected returns.  I.e., not only are their expeceted returns higher, which is consistent with annualized metrics of risk, but their {\it risk} decreases based on the investment horizon.  It's worth looking into further studies of risk metrics over long periods, and coming to terms with the fact that avilable historical data becomes much more restricted when the autocorrelation is relevant over long investment periods.


\section{Brian J. Alleva -- Discount Rate Specification and the Social Security Claiming Decision}

In \cite{Alleva2016}, the author analyzes the Optimal Claiming Age (OCA) for social security benefits, based on a discounted cash flow model across a range of discount rates.  The author notes that previous studies have enforced a government bond-like rate (the rate actually used to fund the social security trust), which may not necessarily hold.  An interesting perspective is note that social security is compulsory, or pre-ordained, vs. something sought out on capital markets.  This is to say that some people claiming social security may not have chosen this investment vehicle if given the option, and so they may not be valuing the income stream strictly as an annuity they went out and purchased.

The trend is stated at the start: discounting by a lower rate (e.g., US bonds) suggests picking at 70, discounting by a higher rate (e.g., an aggressive stock portfolio) suggests picking at 62; i.e. the bounds of social security OCA are set pretty quickly at either end of the discount rate.  Curiously, there is still somewhat elevated risk due to lifetime expectancy associated with deferring benefits until 70; i.e., those using a small discount rate may be under-valuing the probability of not living long enough to enjoy the increased social security payout.  Alternatively, those using a higher discount rate are much more likely to maximize their utility (if that rate reflects reality).  Another curiosity is that the OCA appears to vary rapidly near somewhat conservative discount rates near 4\%, so that the OCA changes from 70 to 62 across a 1\% range.  

\note{[Note]}: to me, the above suggests that the OCA--as defined in the text--could be difficult to determine in practice.  Or, rather, the real claiming age chosen may have more to do with work suitabilty, income opportunity, and/or medicare benefits than a purely discounted cashflow analysis.

\note{[Note]}: it looks like this article is using {\it real} rates of return to discount social security, as the payouts track inflation at least somewhat due to the Cost Of Living Adjustment (COLA).

The plots (Charts 1-3) of death age or probability of optimal claiming age range vs. discount rate are pretty interesting.  For women, there is a much larger range that should seriously consider deferring until 70.  For men, the result are more mixed, skewing toward 62 for somewhat-risky rates of return like 4\%.

The plots (Charts 4 and 5) of OCA and Expected Opportunity Loss (EOL) for men and women seem to indicate that 62 is the OCA unless discounting by an extremely (less than 50/50 stocks/bonds) conservative portfolio.  It would be interesting to consider when this would actually be the case.  I suppose this is the statement that claiming social security at 62 to simple store it in a 1\% interest savings account is not worth it.  Alternatie uses like spending the money instead of taking on debt, or spending social security instead of drawing down medium-to-high risk portfolios, appear to strongly suggest claiming social security sooner than later.

\note{[Note]}: there is some discussion to the above point of how fund usage relates to discount rate in the ``Implied Discount Rates'' section.


\note{[Takeaway]}: I initially had heard that delaying social security made sense, to maximize the payout given that modern lifetimes are expected to increase relative to the benchmarks used to create social security payouts.  After more research in general, but also after reading this paper, I feel like there is a lot of risk related to deferring even if that cash amount may be higher.  Given how moderate of a portfolio is needed to exceed the implied discount rate associated with waiting past 62, it seems like there is little need to delay other than working, which I felt was ignored by this study.  This is to say that if you are planning on retiring, I imagine your retirement accounts will earn sufficiently much so that you should claim immediately.  If you plan on working through your 60s, of course, then I suspect your earnings exceed the tax-penalized payouts from claiming social security early.


\section{Wikipedia article on random walk}

I need to brush up on my basic mathematical statistics/probability, so I thought it would be fun to do a deep dive on random walks, starting from the Wikipedia article \cite{RandomWalk2026}.  I'd drawn some patterns to see that the distribution for a given set of steps is related to the $n$th level of Pascal's triangle \cite{PaslcalsTriangle2026}, but the easier way to understand it is to realize that the final position on an integer line is to consider $n$ independent rolls of $\pm1$.  Then,
\begin{equation}
    \text{E}(S_n) = \sum_{i=1}^n Z_i = 0\quad ,
\end{equation}
and 
\begin{equation}
    \text{E}(S_n^2) = \sum_{i,j = 1}^n Z_i Z_j = \sum_{i=1}^n Z_i^2 = n
\end{equation}
where $S_n$ is the final position relative to starting at 0 on the integer line, and $Z_i$ are independent random variables, whose value is either $\pm1$.  Notably, the fact that each increment of $\pm1$ is independent of prior results is what eliminates any cross terms in the variance.

There is a lot of interesting discussion, but one point I found interesting was the application to the gambler's ruin \cite{GamblersRuin2026}.  Because the variance grows with the number of steps as $\sqrt{n}$, any point is reachable with finite probability for sufficiently large $n$.  This is to say that a gambler playing a fair game (expected outcome of 0) against an opponent with infinite wealth (e.g., a toy model of a casino) will eventually lose everything, as their outcome will randomly walk through 0 no matter how far ahead they start.




\bibliographystyle{unsrturl}
\bibliography{mybib/mybib.bib}

\end{document}
